\chapter{PHƯƠNG PHÁP}

\section{Quy trình phân tích tổng thể}

Quy trình phân tích được thực hiện theo các bước sau:

\begin{enumerate}
    \item \textbf{Thu thập và chẩn đoán dữ liệu}
    \item \textbf{Xử lý dữ liệu (làm sạch, chuẩn hóa)}
    \item \textbf{Phân tích khám phá (EDA)}
    \item \textbf{Giảm chiều dữ liệu (PCA)}
    \item \textbf{Xác định K tối ưu (Elbow + Silhouette)}
    \item \textbf{Phân cụm K-means}
    \item \textbf{Xác thực kết quả (Hierarchical, DBSCAN)}
    \item \textbf{Phân tích đặc trưng cụm}
    \item \textbf{Phát hiện dị biệt (Outliers)}
    \item \textbf{Đưa ra khuyến nghị}
\end{enumerate}

\section{Dữ liệu sử dụng}

\subsection{Nguồn dữ liệu}

Dữ liệu sử dụng là CC GENERAL.csv - bộ dữ liệu công khai về khách hàng thẻ tín dụng.

\subsection{Mô tả dữ liệu}

\begin{table}[H]
\centering
\caption{Danh sách các biến trong bộ dữ liệu}
\small
\begin{tabular}{|p{4.5cm}|p{2cm}|p{6cm}|p{1.5cm}|}
\hline
\textbf{Biến} & \textbf{Loại} & \textbf{Mô tả} & \textbf{Đơn vị} \\
\hline
CUST\_ID & Chuỗi & Mã khách hàng & ID \\
\hline
BALANCE & Số thực & Số dư trong tài khoản & USD \\
\hline
BALANCE\_FREQUENCY & Số thực & Tần suất có số dư & \% \\
\hline
PURCHASES & Số thực & Tổng giá trị mua hàng & USD \\
\hline
ONEOFF\_PURCHASES & Số thực & Tổng tiền mua một lần & USD \\
\hline
INSTALLMENTS\_PURCHASES & Số thực & Tổng tiền mua trả góp & USD \\
\hline
CASH\_ADVANCE & Số thực & Tổng tiền rút tiền mặt (cash advance) & USD \\
\hline
PURCHASES\_FREQUENCY & Số thực & Tần suất mua hàng & \% \\
\hline
CASH\_ADVANCE\_FREQUENCY & Số thực & Tần suất rút tiền & \% \\
\hline
CASH\_ADVANCE\_TRX & Số nguyên & Số lần rút tiền & Lần \\
\hline
PURCHASES\_TRX & Số nguyên & Số lần mua hàng & Lần \\
\hline
CREDIT\_LIMIT & Số thực & Giới hạn tín dụng & USD \\
\hline
PAYMENTS & Số thực & Tổng tiền thanh toán & USD \\
\hline
MINIMUM\_PAYMENTS & Số thực & Thanh toán tối thiểu & USD \\
\hline
PRC\_FULL\_PAYMENT & Số thực & Tỷ lệ thanh toán toàn bộ & \% \\
\hline
TENURE & Số nguyên & Số tháng làm khách hàng & Tháng \\
\hline
\end{tabular}
\label{tab:variables}
\end{table}

\textbf{Kích thước dữ liệu:}
\begin{itemize}
    \item Số hàng ban đầu: 8,950 khách hàng (sau xử lý NaN)
    \item Số cột: 18 (17 biến + 1 ID)
    \item Nguồn: Kaggle - Credit Card Dataset for Clustering
\end{itemize}

\section{Kỹ thuật xử lý dữ liệu}

\subsection{Kiểm tra giá trị thiếu (Missing Values)}

Các giá trị thiếu được xác định và xử lý như sau:
\begin{itemize}
    \item \textbf{MINIMUM\_PAYMENTS:} Điền bằng median
    \item \textbf{CREDIT\_LIMIT:} Điền bằng median
    \item Các biến khác: Không có giá trị thiếu
\end{itemize}

Công thức điền median:
\begin{equation}
x_i = \text{median}(x) \text{ nếu } x_i \text{ là missing}
\end{equation}

\subsection{Phép biến đổi logarit (Log Transformation)}

Do dữ liệu tài chính thường lệch phải (right-skewed), áp dụng biến đổi logarit:

\begin{equation}
x'_i = \log(1 + x_i)
\end{equation}

Cộng 1 để tránh $\log(0)$.

\textbf{Lợi ích:}
\begin{itemize}
    \item Giảm độ lệch của phân phối
    \item Giúp các giá trị lớn và nhỏ có trọng lượng tương đương
    \item Cải thiện hiệu suất mô hình
\end{itemize}

\subsection{Chuẩn hóa dữ liệu (Standardization)}

Sử dụng StandardScaler để chuẩn hóa:

\begin{equation}
z = \frac{x - \mu}{\sigma}
\end{equation}

Trong đó:
\begin{itemize}
    \item $\mu$ = giá trị trung bình của biến
    \item $\sigma$ = độ lệch chuẩn của biến
\end{itemize}

\textbf{Lý do:}
\begin{itemize}
    \item K-means nhạy cảm với thang đo biến
    \item Chuẩn hóa không làm thay đổi hình dạng phân phối
    \item Đảm bảo các biến có cùng ảnh hưởng
\end{itemize}

\section{Giảm chiều dữ liệu}

\subsection{Tại sao cần giảm chiều?}

\begin{itemize}
    \item Dữ liệu gốc có 17 chiều, khó trực quan hóa
    \item Giảm thời gian tính toán
    \item Giảm nhiễu và overfitting
    \item Giúp phát hiện cấu trúc ẩn
\end{itemize}

\subsection{Áp dụng PCA}

PCA được áp dụng để giảm chiều dữ liệu từ 17 chiều xuống 2 chiều:
\begin{itemize}
    \item \textbf{Giảm chiều:} Giảm từ 17 biến xuống 2 thành phần chính ($PC1, PC2$)
    \item \textbf{Giữ lại thông tin:} 2 PC giữ lại khoảng 56.46\% phương sai của dữ liệu gốc
    \item \textbf{Trực quan hóa:} Dễ dàng vẽ biểu đồ phân tán 2D
    \item \textbf{Phân cụm:} K-means được thực hiện trên không gian PCA 2 chiều này
\end{itemize}

Việc sử dụng PCA 2D cho cả trực quan hóa và clustering giúp:
\begin{itemize}
    \item Giảm nhiễu từ các biến ít quan trọng
    \item Tăng tốc độ tính toán
    \item Giảm overfitting
    \item Dễ dàng diễn giải và trình bày kết quả
\end{itemize}

\begin{equation}
\text{Variance Explained by 2 PCs} = \frac{\lambda_1 + \lambda_2}{\sum_{i=1}^{17} \lambda_i}
\end{equation}

\section{Xác định số cụm tối ưu}

\subsection{Phương pháp Elbow}

Tính inertia cho mỗi K từ 2 đến 10:

\begin{equation}
\text{Inertia}(K) = \sum_{k=1}^{K} \sum_{x \in C_k} \|x - \mu_k\|^2
\end{equation}

Điểm "gãy khúc" nơi tốc độ giảm chậm lại chỉ ra K tối ưu.

\subsection{Chỉ số Silhouette}

Tính Silhouette Score trung bình cho mỗi K:

\begin{equation}
\text{SilhouetteScore}(K) = \frac{1}{n} \sum_{i=1}^{n} s(i)
\end{equation}

K tối ưu thường có Silhouette Score cao nhất.

\section{Phân cụm K-means}

\subsection{Thuật toán K-means}

\begin{algorithm}
\textbf{Input:} Dữ liệu $X = \{x_1, \ldots, x_n\}$, số cụm $K$
\newline
\textbf{Output:} Nhãn cụm cho mỗi điểm
\begin{enumerate}
    \item \textbf{Khởi tạo:} Chọn ngẫu nhiên K tâm cụm $\mu_1, \ldots, \mu_K$
    \item \textbf{Lặp:}
    \begin{enumerate}
        \item \textbf{Gán:} Gán mỗi $x_i$ tới cụm gần nhất: $c_i = \arg\min_k \|x_i - \mu_k\|$
        \item \textbf{Cập nhật:} $\mu_k = \frac{1}{|C_k|} \sum_{x \in C_k} x$
    \end{enumerate}
    \item \textbf{Dừng:} Khi tâm cụm hội tụ hoặc số lần lặp đạt giới hạn
\end{enumerate}
\end{algorithm}

\subsection{Khởi tạo K-means}

Sử dụng k-means++ để khởi tạo tâm cụm:
\begin{itemize}
    \item Chọn tâm đầu tiên ngẫu nhiên
    \item Mỗi tâm tiếp theo: chọn điểm xa tâm gần nhất
    \item Giảm nguy cơ hội tụ cục bộ không tốt
\end{itemize}

\section{Xác thực kết quả phân cụm}

\subsection{So sánh với phân cụm phân cấp}

Áp dụng phân cụm phân cấp với Ward linkage để xác nhận kết quả K-means.

\subsection{Phát hiện dị biệt bằng DBSCAN}

Sử dụng DBSCAN với tham số:
\begin{itemize}
    \item $\epsilon = 0.5$ (bán kính láng giềng)
    \item $\text{min\_samples} = 10$ (số điểm tối thiểu)
\end{itemize}

Mục đích:
\begin{itemize}
    \item Phát hiện điểm bất thường tự động
    \item Không cần xác định K trước
    \item Xác định khách hàng có hành vi khác biệt đáng kể
\end{itemize}

\section{Phân tích và đặc trưng cụm}

Đối với mỗi cụm, tính toán:
\begin{itemize}
    \item \textbf{Thống kê mô tả} (mean, std, min, max) cho mỗi biến
    \item \textbf{Kích thước cụm} = số khách hàng trong cụm
    \item \textbf{Đặc trưng chính} = biến phân biệt nhất cụm này với cụm khác
    \item \textbf{Tên/Label} cho mỗi cụm dựa trên hành vi
\end{itemize}

\section{Công cụ và môi trường}

\textbf{Ngôn ngữ lập trình:} Python 3.8+

\textbf{Thư viện chủ yếu:}
\begin{itemize}
    \item \textbf{pandas:} Xử lý dữ liệu bảng
    \item \textbf{numpy:} Tính toán số học
    \item \textbf{scikit-learn:} Thuật toán ML (KMeans, PCA, DBSCAN)
    \item \textbf{scipy:} Tính toán thống kê (linkage, dendrogram)
    \item \textbf{matplotlib, seaborn:} Trực quan hóa dữ liệu
\end{itemize}

\textbf{Môi trường thực thi:} 
\begin{itemize}
    \item Jupyter Notebook
    \item Python Virtual Environment (.venv)
    \item CPU đa lõi (xử lý song song)
\end{itemize}
